% ================================================================
%  conteudo/introducao.tex
% ================================================================

\section{Introdução}

A introdução é a primeira seção textual do relatório científico e tem como objetivo contextualizar o leitor sobre o tema da prática realizada. Ela deve apresentar a relevância do assunto, situando-o no campo científico ou tecnológico pertinente, e fornecer informações suficientes para que o leitor compreenda o porquê do experimento ser realizado.

Um bom parágrafo introdutório começa com uma afirmação ampla sobre a área de estudo e, progressivamente, afunila o raciocínio até chegar ao tema específico da prática. Por exemplo: os processos fermentativos constituem uma das ferramentas biotecnológicas mais antigas e relevantes da humanidade, sendo amplamente utilizados na produção de alimentos, bebidas, fármacos e biocombustíveis (Silva; Aquino, 2018). Sua aplicação industrial fundamenta-se na capacidade de microrganismos em converter substratos orgânicos em produtos de interesse por meio de reações enzimáticas controladas (Schmidell \textit{et al.}, 2001).

A fermentação pode ocorrer em condições aeróbias ou anaeróbias, dependendo do microrganismo envolvido e do produto desejado. Segundo Lehninger, Nelson e Cox (2014), as vias metabólicas envolvidas determinam tanto a eficiência do processo quanto os subprodutos gerados. O controle de variáveis como temperatura, pH, concentração de substrato e aeração é essencial para a otimização dos rendimentos fermentativos (Costa \textit{et al.}, 2020).

Nesse contexto, a prática realizada teve como pano de fundo a investigação de parâmetros cinéticos de crescimento microbiano em reator de bancada, buscando compreender a influência das condições operacionais sobre a produtividade do processo. O entendimento desses fundamentos é indispensável para a formação do engenheiro de alimentos ou bioquímico que atuará no desenvolvimento e controle de bioprocessos industriais.

